% ACL Rolling Review submission — official ACL Author Kit.
% \usepackage[review]{acl} for anonymous review (current).
% Switch to \usepackage{acl} for the camera-ready.
\documentclass[11pt]{article}
\PassOptionsToPackage{hyphens}{url}   % let long URLs break at hyphens
\usepackage[review]{acl}

\usepackage[T1]{fontenc}
\usepackage[utf8]{inputenc}
\usepackage[english]{babel}
\usepackage{times}
\usepackage{latexsym}
\usepackage{amsmath,amssymb}
\usepackage{booktabs}
\usepackage{graphicx}
\usepackage{multirow}
\usepackage{microtype}

\title{Sentence Classification Is the Wrong Task:\\
Query-Conditioned Verification of Literature-Mined\\Biotic-Interaction Triples}

% Anonymised automatically by acl.sty under the [review] option.
\author{
  Anonymous ACL submission
}

\begin{document}
\maketitle

\begin{abstract}
Literature-mining systems that populate biotic-interaction databases retrieve candidate
(species, relation, species) triples with supporting passages, then filter them with a
sentence-level interaction classifier before curation. We show that this filter answers the
wrong question. On 100 expert-graded triples retrieved by a deployed system, the filter accepts
91 and reaches precision 0.714, which is not significantly above accepting every candidate
(0.650; one-sided exact binomial $p = 0.119$); it rejects 9 of 35 unsupported triples. We
decompose those 35 errors into entity-linking (10), predicate (9) and \emph{pair-binding} (16)
failures, and show that a per-component oracle is bounded by $65r/(65r+16)$ at recall $r$, because
16 errors have every component individually correct and are therefore invisible to it at any
operating point. The filter's decisions track whole-triple support (Matthews $0.429$) more closely
than the sentence-level property it was trained on ($0.128$). Conditioning the decision on the queried triple removes most of the gap without any
training: a locally hosted 32B model, prompted zero-shot with the candidate triple, rejects
32 of 35 at precision 0.949 ($p < 10^{-4}$), and reaches F1 0.886 (95\% CI [0.846, 0.921]) on a
299-item pair-level benchmark where the best of our supervised models reaches 0.823. We then
explain why \emph{supervised} pair conditioning has repeatedly failed in this setting: our
paired training corpora contain 606 and 11{,}602 distinct taxon pairs against 79{,}460 in the
retrieval pool, and the effect of conditioning flips sign with that diversity. We release the
graded benchmark, the audited corpora, and a distilled cross-encoder.
\end{abstract}

\section{Introduction}
\label{sec:intro}

Biotic-interaction databases such as GloBI \citep{poelen2014global} are populated in part by
text mining: a retrieval system proposes candidate triples $(s_1, r, s_2)$ together with the
passage that allegedly supports them, and a filter decides which candidates are worth a
curator's attention. The filter is almost always a \emph{sentence-level} interaction
classifier: it is shown the passage and asked whether it describes a biotic interaction.

This paper argues that the sentence-level formulation is a category error for this task, and
supports the argument with measurements on a deployed system. Our contributions:

\begin{itemize}
\item \textbf{The deployed filter does not filter.} On its own expert-graded evaluation set it
      accepts 91 of 100 candidates at precision 0.714, statistically indistinguishable from
      accepting all of them (\S\ref{sec:nofilter}).
\item \textbf{A failure decomposition with a hard ceiling.} Of 35 unsupported triples, 16 have
      correct entities \emph{and} a correct relation; the triple is wrong only in how the
      arguments bind. Any component-wise verifier is capped at precision 0.802
      (\S\ref{sec:decomp}).
\item \textbf{Query conditioning closes most of the gap, zero-shot.} Supplying the candidate
      triple in the prompt raises rejection from 9/35 to 32/35 and precision to 0.949, and
      beats every model we trained on a 299-item benchmark (\S\ref{sec:querycond}).
\item \textbf{Why supervised pair conditioning failed.} The effect of pair conditioning flips
      sign with taxon-pair diversity of the training corpus; at 606 distinct pairs the pair slot
      is noise, at 11{,}602 it helps by $+0.098$ AUPRC (\S\ref{sec:diversity}).
\end{itemize}

\section{Setting}
\label{sec:setting}

\paragraph{The pipeline under audit.} BiotXplorer \citep{ruch2024biotxplorer} proposes candidate
biotic-interaction triples from the literature indexed by SIB Literature Services
\citep{gobeill2020sibils}. For each passage it matches taxon surface forms against a taxonomy and
interaction surface forms against a vocabulary derived from the Relation Ontology, and emits every
resulting $(s_1, r, s_2)$ combination together with the passage, the three matched surface
strings, their canonical forms, and a heuristic passage score. Candidates are therefore proposed
by \emph{co-occurrence}: a passage naming two taxa near an interaction term produces a candidate
whether or not it asserts that interaction between those two. A filter then decides what reaches
a curator, and in the deployed system that filter is a sentence-level interaction classifier.

\paragraph{\textsc{Biotx100}.} One hundred retrieved candidates graded by a domain expert on four
axes: whether each taxon surface form was correctly resolved ($92/100$ and $98/100$), whether the
relation was correct ($90/100$), and whether the whole triple is supported by its passage
($65/100$). It also records the deployed filter's own accept/reject decision, which is what makes
the comparison in \S\ref{sec:nofilter} possible. Every row has $\mathrm{rank}=1$, so this is
precision at one retrieved candidate, not a ranked evaluation.

\paragraph{\textsc{Test299}.} A larger pair-level set assembled from expert curation only. We
recovered the taxon pair for 299 of 500 items in an existing multi-source test set by joining
back to the curation files they were drawn from; the recovered subset is exactly the
human-curated portion, excluding 197 LLM-generated items. The relation term is recoverable for
287 of those, and we report on that subset wherever a query is required. It spans three sources
that differ sharply in prevalence, from $0.31$ to $0.84$ positive.

\paragraph{Evaluation protocol.} Decision thresholds are fixed before evaluation and never
selected on the set being reported: they come from a held-out development split or from a
closed-form prior-shift rule $\tau = \sigma(-[\mathrm{logit}(\pi_{\mathrm{target}}) -
\mathrm{logit}(\pi_{\mathrm{train}})])$. All comparisons report the trivial majority-class
baseline, all intervals are 10{,}000-sample bootstraps, and all model figures are means over
three seeds. Reporting the trivial baseline is not a formality here: on one inherited source an
all-positive predictor scores F1 $0.914$, above every system we tested.

\section{The deployed filter does not filter}
\label{sec:nofilter}

\begin{table}[t]\centering\small
\begin{tabular}{lrrrr}
\toprule
System & Acc. & Prec. & Rej. & $p$ \\
\midrule
Accept everything      & 100 & 0.650 & 0/35 & --- \\
Deployed classifier    &  91 & 0.714 & 9/35 & 0.119 \\
\midrule
LLM, pair only         &  71 & 0.817 & 22/35 & 0.002 \\
LLM, full triple       &  59 & \textbf{0.949} & \textbf{32/35} & $<10^{-4}$ \\
\bottomrule
\end{tabular}
\caption{\textsc{Biotx100}. ``Acc.'' is candidates accepted, ``Rej.'' unsupported triples
rejected. $p$ is a one-sided exact binomial test against the accept-everything precision of
0.650. The deployed filter is not significantly better than no filter at all.}
\label{tab:nofilter}
\end{table}

Table~\ref{tab:nofilter} gives the central observation. The deployed classifier accepts 91 of
100 candidates. Its precision, 0.714, is higher than the 0.650 obtained by accepting
everything, but not significantly so ($p = 0.119$). It recovers 9 of the 35 unsupported
triples. Its recall on supported triples is 1.000 --- it rejects nothing that should be kept,
and almost nothing that should not.

This is not a threshold-tuning artifact. The classifier's decisions correlate more strongly with
axes it was never trained on than with the one it was. Against the sentence axis it was trained
on --- does this passage describe a biotic interaction --- its decisions reach a Matthews
correlation of only $0.128$; against whether the two taxa were correctly resolved, $0.245$; and
against whole-triple support, $0.429$.

The sharpest form of the observation restricts attention to the 90 passages that \emph{do}
describe a biotic interaction, so the sentence-level question is answered affirmatively for all
of them. Among these the filter accepts 83 of 90, and of the 25 that nonetheless fail to support
their candidate triple it rejects 7. Once the sentence-level question is settled, the filter has
almost no remaining discriminative power over the question the pipeline actually asks.

\section{What the errors are made of}
\label{sec:decomp}

\begin{table}[t]\centering\small
\begin{tabular}{lrl}
\toprule
Failure mode & $n$ & Detectable by \\
\midrule
Entity linking ($s_1$ or $s_2$ wrong) & 10 & taxon resolver \\
Predicate (relation wrong)            &  9 & relation classifier \\
\textbf{Pair binding}                 & \textbf{16} & \textbf{joint judgment only} \\
\midrule
Total unsupported                     & 35 & \\
\bottomrule
\end{tabular}
\caption{Decomposition of the 35 unsupported triples in \textsc{Biotx100}. The 16
pair-binding failures have both taxa correctly resolved \emph{and} the correct relation type;
the passage simply does not assert that relation \emph{between those two}.}
\label{tab:decomp}
\end{table}

Table~\ref{tab:decomp} decomposes the failures. The largest class, 16 of 35, consists of
triples in which every component is individually correct. Typically the passage asserts the
relation between one of the queried taxa and some third organism, or lists several taxa of
which only a subset participate.

This bounds any component-wise architecture at every operating point. A verifier that checks
$s_1$, $r$ and $s_2$ separately --- however good each check is --- rejects at most the 19
component-detectable errors, and has no signal at all that separates the 16 pair-binding failures
from supported triples. Accepting a fraction $r$ of the 65 supported triples it therefore also
accepts all 16, giving precision $65r/(65r+16)$: $0.802$ at full recall, $0.754$ at $r = 0.754$.
The remaining headroom is reachable only by a decision that considers the triple jointly.
Figure~\ref{fig:bound} plots that bound. The deployed classifier sits \emph{below} it: at full
recall it reaches 0.714 where a per-component oracle would reach 0.802, so it is not merely
failing to exploit joint information, it is losing information the components alone would supply.

\begin{figure}[t]\centering
\includegraphics[width=\columnwidth]{fig/fig1_oracle_bound.pdf}
\caption{Precision against recall of supported triples on \textsc{Biotx100}. The dashed curve is
the per-component oracle bound $65r/(65r+16)$: a verifier with perfect knowledge of both taxa and
the relation still admits all 16 pair-binding failures, at every operating point. The deployed
classifier falls below it; the query-conditioned verifier clears it by $0.21$ at matched recall.}
\label{fig:bound}
\end{figure}

\section{Query conditioning}
\label{sec:querycond}

\begin{table}[t]\centering\small
\begin{tabular}{lrr}
\toprule
System & F1 & 95\% CI \\
\midrule
Trivial (all positive)          & 0.706 & --- \\
Prior multi-task model          & 0.791 & [0.737, 0.840] \\
Template-trained BiomedBERT     & 0.821 & [0.772, 0.865] \\
Best supervised (this work)     & 0.823 & --- \\
\midrule
LLM, query-conditioned, 0-shot  & \textbf{0.886} & [0.846, 0.921] \\
Distilled cross-encoder (110M)  & 0.781 & [0.729, 0.828] \\
\bottomrule
\end{tabular}
\caption{\textsc{Test299}. Thresholds fixed a priori. The zero-shot query-conditioned model's
interval excludes every supervised system.}
\label{tab:test299}
\end{table}

We supply the candidate triple in the prompt and ask a locally hosted 32B instruction model,
zero-shot at temperature 0, whether the passage supports it. No training, no in-context
examples, no access to the benchmark.

On \textsc{Biotx100} this raises rejection of unsupported triples from 9/35 to 32/35 and
precision from 0.714 to 0.949 (Table~\ref{tab:nofilter}). The gain is not merely from naming
the two taxa: a prompt giving only the pair, without the relation, reaches 0.817 --- so
conditioning on the \emph{relation} is worth a further $+0.13$ precision.

On \textsc{Test299} (Table~\ref{tab:test299}) the same model reaches F1 0.886, with a bootstrap
interval that excludes every supervised system we trained, including a multi-task architecture
developed over several corpus generations.

\section{Conditioning the input is not enough}
\label{sec:supervision}

Given \S\ref{sec:querycond}, the obvious move is to train a model with the pair in its input.
We did, and by itself it does nothing.

We added the taxon pair to the multi-task classifier as a genuine segment pair --- the pair as
segment~A, the passage as segment~B, entity-recognition supervision masked over the query --- and
changed nothing else: same corpus, same architecture, same labels, three seeds. AUPRC moves from
$0.848$ to $0.844$ (Figure~\ref{fig:supervision}, left pair), against a seed standard deviation
of $0.010$. The model is given exactly the information it was missing and does not use it.

\begin{figure}[t]\centering
\includegraphics[width=\columnwidth]{fig/fig2_supervision.pdf}
\caption{Only the left pair is a controlled comparison: one factor changed, same corpus,
architecture and labels. The query-conditioned verifier differs additionally in supervision,
corpus and architecture, and is shown for scale rather than as a matched cell.}
\label{fig:supervision}
\end{figure}

The reason is that we changed the question the model is \emph{asked} without changing the
question its labels \emph{answer}. Those labels come from an ensemble scoring sentences, so the
target is unchanged when the queried pair changes: two training items with the same passage and
different pairs carry the same label. Gradient descent has no reason to attend to segment~A,
because segment~A never moves the target. The verifier of \S\ref{sec:querycond} works because its
teacher was itself asked the pair-specific question, so its labels vary with the query.

This also accounts for a negative result recorded earlier in this project's own history, that
prepending entity names to the input degrades performance. Under a sentence-level target the
prepended string is uninformative tokens competing for capacity, which is what a degradation
looks like.

Corpus diversity compounds it. Conditioning helps where the paired corpus is diverse enough for
pair identity to generalise, and hurts where it is not: $-0.003$ AUPRC on a corpus with 606
distinct taxon pairs, $+0.098$ on one with $11{,}602$, against $79{,}460$ in the retrieval pool.
At 606 pairs the pair slot is roughly ten tokens of noise; at $11{,}602$ it is a feature. Both
conditions share the sentence-level target, so diversity governs how much a pair-conditioned
input can help, and the supervision governs whether it can help at all.

\paragraph{What we did not run.} The fourth cell --- sentence-level input with pair-level
supervision --- would separate the two factors completely. It requires re-labelling a
sentence-only corpus with the query-conditioned teacher, and we report the three cells we have.

\section{Distillation}
\label{sec:distill}

We label that sample with the query-conditioned teacher (28.9\% positive, no failures) and train
a 110M-parameter cross-encoder reading the candidate triple as segment~A and the passage as
segment~B, with a split grouped on taxon pair so no pair spans training and development.

\begin{table}[t]\centering\small
\begin{tabular}{lrrr}
\toprule
System & Prec. & Rej. & AUPRC \\
\midrule
Accept everything          & 0.650 & 0/35 & --- \\
Deployed classifier        & 0.714 & 9/35 & --- \\
Per-component oracle       & 0.754 & 19/35 & --- \\
\midrule
LLM teacher (32B)          & 0.949 & 32/35 & --- \\
Student (110M)             & \textbf{0.967} & \textbf{33.3/35} & \textbf{0.971} \\
\bottomrule
\end{tabular}
\caption{\textsc{Biotx100}. Student figures are means over three seeds (precision sd $0.022$),
using the canonical taxon and relation strings; the oracle row is evaluated at the student's
matched recall of $0.754$ rather than at full recall. The student beats the matched-recall
component oracle by $0.213$ at roughly $290\times$ fewer parameters than its teacher.}
\label{tab:student}
\end{table}

Table~\ref{tab:student} gives the deployed-task result. The student reaches precision $0.967$ and
rejects $33.3$ of $35$ unsupported triples. Two things are worth separating. Compared at its own recall of $0.754$, where a per-component
oracle reaches only $0.754$, it gains $0.213$ precision --- the paper's central prediction, and
not reachable by checking entities and relation separately. And it does so at about $290\times$
fewer parameters than the teacher. Precision depends on how the query is written: supplying the
canonical taxon and relation strings gives $0.967$ at recall $0.754$, while the retrieved surface
forms give $0.928$ at recall $0.851$ --- a higher-recall, slightly lower-precision point.

On \textsc{Test299} the student reaches F1 $0.781$ and AUPRC $0.905$: better ranking than any
supervised model we trained, but lower F1 than the template baseline, because its operating point
is set for the $28.9\%$ prior of the retrieval pool rather than that benchmark's $54.4\%$.

\begin{table}[t]\centering\small
\begin{tabular}{lrrrrr}
\toprule
Source & $n$ & $\pi$ & Trivial & Student & AUPRC \\
\midrule
EP-A                  & 99 & 0.47 & 0.644 & \textbf{0.842} & 0.917 \\
EP-passage            & 95 & 0.84 & \textbf{0.914} & 0.819 & 0.936 \\
BioTx-random          & 93 & 0.31 & 0.475 & 0.622 & 0.818 \\
\midrule
Overall               & 287 & 0.54 & 0.704 & 0.796 & 0.908 \\
\bottomrule
\end{tabular}
\caption{Per-source results on \textsc{Test299}, student mean of three seeds at the shared
prior-shift threshold $\tau = 0.254$. Bold marks the higher of student and trivial baseline. AUPRC
is threshold-free and stays above $0.81$ on every source; F1 does not, and the reason is
calibration rather than ranking.}
\label{tab:persource}
\end{table}

\paragraph{Per-source behaviour, and what it says about deployment.}
Table~\ref{tab:persource} breaks the result down by source. Two things are visible and neither is
flattering. On EP-passage, which is $84\%$ positive, an all-positive predictor scores $0.914$ and
every system we tested --- ours included --- scores below it; that source cannot discriminate
between methods and we report it only so that its inclusion in the aggregate is visible. On
BioTx-random the student scores $0.622$ where a plain template-trained baseline reaches $0.929$,
and this is not a contamination artifact: exact and species-masked overlap between that source and
the baseline's training corpus is $0/93$ in both cases.

The mechanism is calibration, not discrimination. At the shared threshold the student is
high-precision and low-recall on BioTx-random ($P = 0.875$, $R = 0.483$), and its scores there are
compressed: mean confidence on supported triples is $0.435$ against $0.774$ on EP-A. Its AUPRC on
that source is $0.818$, and a per-source threshold would recover F1 $0.784$ from $0.622$ --- so
more than half the gap is where the boundary sits, not how the items are ranked. A single global
operating point is the wrong deployment contract for a system whose input distribution varies this
much; the practical consequence is that a curator-facing filter should expose a calibrated score
with an adjustable threshold rather than a fixed verdict.

\paragraph{The query is doing the work.} Supplying the candidate relation matters as much as
supplying the pair. Scored on the same 287 \textsc{Test299} items with a generic relation string
substituted for the retrieved one, the student falls from F1 $0.781$ to $0.532$: a swing of
$0.249$ from query specificity alone, with passage and taxon pair held fixed.

\section{Related work}
\label{sec:related}

\paragraph{Filters that answer an easier question.}
That a classifier can score well by exploiting how its data were assembled is among the
best-established results in NLP. Partial-input baselines recover much of the label from the
hypothesis alone in natural language inference \citep{gururangan-etal-2018-annotation,poliak-etal-2018-hypothesis} and from the passage alone in reading comprehension
\citep{kaushik-lipton-2018-much}, and such rules hold in distribution while failing under
deployment conditions \citep{geirhos-etal-2020-shortcut,du-etal-2024-shortcut}. The
closest task-level precedent is fact verification, where a claim-only classifier matches
evidence-aware models because of how the corpus was built \citep{schuster-etal-2019-towards}.
Section~\ref{sec:nofilter} reports that phenomenon rather than a new one: a filter
statistically indistinguishable from its majority class is a partial-input baseline that
reached production. What is new is the identity of the missing input --- the query the pipeline
itself issued --- and that the cue, where present, originates in the curation protocol rather
than in an annotator \citep{wiegand-etal-2019-detection,parmar-etal-2023-dont}. We read our lexicon
diagnostic in that spirit only: a keyword baseline that fails certifies nothing
\citep{feng-etal-2019-misleading}.

\paragraph{Relation classification never assumed the sentence was enough.}
Distant supervision assumes that a sentence containing two related entities expresses the
relation \citep{mintz2009distant}, and its standard relaxation exists precisely because it does
not, separating ``these two are related'' from ``this sentence says so''
\citep{riedel2010modeling}. Since TACRED an instance has been the triple (sentence, $e_1$,
$e_2$) rather than the sentence \citep{zhang2017position}, and how the pair is injected into the
encoder is worth several F1 points \citep{soares2019matching,wu2019enriching,zhou2022improved}.
\citet{rosenman2020exposing} name our failure exactly, as an ``event heuristic'' that
classifies on (sentence, relation) while ignoring the arguments, and show that QA-style models
repair it on a challenge set. We claim no part of that diagnosis. Our contribution is that the
heuristic is observable in a deployed system's own expert-graded output rather than in an
adversarially constructed set; that its cost is quantifiable as a per-component oracle ceiling
(\S\ref{sec:decomp}); and that the entity-type shortcut usually invoked to explain such errors
\citep{peng2020learning,mtumbuka2024entity} is unavailable here, because both arguments are
taxa and type therefore carries no discriminative signal.

\paragraph{Conditioning the decision on the query.}
Casting a relation as a question about a named entity \citep{levy2017zero}, or as an entailment
hypothesis verbalising the relation between two of them
\citep{obamuyide2018zero,sainz2021label}, is the reformulation our prompt instantiates; the
resulting decision --- does this passage support this specific claim --- has the shape of
scientific claim verification \citep{wadden2020fact}. Whether a prompted model beats a
supervised one at biomedical extraction is contested: it is reported for relation extraction
generally \citep{wadhwa2023revisiting,li2023zeroshotre} and denied for biomedical information
extraction \citep{jimenezgutierrez2022thinking,brokman2025benchmark}. That disagreement concerns
open-schema label assignment, whereas we ask a closed binary question about a supplied triple,
which is why we treat the reformulation rather than model scale as the operative variable.
Distilling such a verifier into a cross-encoder \citep{nogueira2019passage} follows established
practice for turning a large model's judgments into a small student
\citep{west2022symbolic,gekhman2023trueteacher}.

\paragraph{Text mining for interaction databases.}
GloBI \citep{poelen2014global} is populated in part by text mining, through semantic-publishing
workflows that route mined interactions into the database \citep{dimitrova2020semantic}; the
system audited here \citep{ruch2024biotxplorer} retrieves over a biodiversity-indexed MEDLINE
\citep{gobeill2020sibils}. Judging whether an extracted fact is correct before it reaches a
curator is an old idea \citep{rodriguezesteban2006imitating}, as is evaluating such tools by
whether they help the curator \citep{hirschman2012biocuration}; and independently of us, \citet{cuzick2023interspecies} argue that the unit of interspecies
curation is the bound pair rather than the passage. Large language models are now applied to
interaction extraction at scale \citep{keck2025extracting,farrell2024landscape}. We are
nonetheless not aware of published work that verifies literature-mined biotic-interaction
triples against their source passage before ingestion, which is the gap this paper occupies.

\section{Conclusion}
\label{sec:conclusion}
The filter in a deployed biotic-interaction mining pipeline is not significantly better than no
filter, because it is asked whether a passage describes an interaction rather than whether it
supports the candidate triple. Sixteen of 35 failures are invisible to any component-wise check.
Conditioning on the query recovers most of the gap with no training at all.

\section*{Limitations}

\paragraph{Benchmark size.} \textsc{Biotx100} contains 100 graded candidates, 35 of them
unsupported, and \textsc{Test299} contains 299 items. On \textsc{Biotx100} one item moves
precision by about a point. The non-significance of the deployed filter's advantage
($p = 0.119$) is a statement about this sample, not evidence that the filter has no effect; a
larger graded set could resolve it either way.

\paragraph{One annotator, no agreement statistic.} All gold labels on both sets come from a
single domain expert; no item was graded twice, so we report no inter-annotator agreement and
cannot bound label noise. This bears hardest on the decomposition: the $16/19$ split between
pair-binding and component-detectable failures is one person's judgment about \emph{where} a
triple fails, and that split fixes the $0.802$ ceiling.

\paragraph{One system, one domain.} We audit one retrieval system, one interaction vocabulary
and one domain, taxon--taxon interactions in MEDLINE. We cannot show that filters over other pools fail the same way,
although the mechanism --- a pair-agnostic decision over passages mentioning more than one
candidate pair --- is not specific to ours.

\paragraph{Pretraining contamination.} The passages are published MEDLINE text and the verifier
is a pretrained 32B model. We cannot exclude that it saw them, or the interactions they report,
during pretraining. Its precision should therefore be read as an optimistic estimate of its
performance on genuinely new literature.

\paragraph{Teacher dependence.} The distilled student inherits the teacher's decision boundary
including its errors, which are systematic rather than random; training encoders on
model-generated labels can degrade and destabilise them \citep{lu2025feeding}. We report the
student's benchmark score, not its agreement with an expert on fresh retrieval output.

\paragraph{Our supervised baselines are weak in a specific way.} Our training corpora carry a
lexical shortcut installed by a filtering step applied to positives only, so our supervised
numbers are a lower bound on what a query-conditioned model trained on cleanly built data could
reach: \S\ref{sec:querycond} compares against the models this project produced, not the best
achievable ones. Rebuilding the corpus is a known remedy \citep{lebras-etal-2020-adversarial},
but does not reliably remove the corresponding bias from the model
\citep{serrano-etal-2023-stubborn}.

\bibliographystyle{acl_natbib}
\bibliography{references}
\end{document}
